\documentclass[12pt]{article}
\usepackage[utf8]{inputenc}
\usepackage[T1]{fontenc}
\usepackage{hyperref}
\title{Ughele Language Resources for K-12 Curriculum Development}
\author{Compiled by Claude Sonnet 5 (claude-sonnet-5) for Andr\'as Kornai}
\date{}
\begin{document}
\maketitle
\section{Introduction}
Ughele (also Ugele) is an Oceanic Austronesian language of northern Rendova Island, Western Province, Solomon Islands. Sources give about 1,200 speakers, a figure dating from 1999. The language has no dialect variation. Before Benedicte Haraldstad Frostad's 2012 grammar it had no writing system and had not been documented.
\section{Primary Educational Resources}
No primers or school materials were found.
\section{Linguistic and Reference Materials}
\begin{itemize}
\item Wikipedia article on the language, with a phonology section based on Frostad's grammar: \url{https://en.wikipedia.org/wiki/Ughele_language}
\item Omniglot page on Ughele: \url{https://www.omniglot.com/writing/ughele.htm}
\item Omniglot page on Ughele numbers: \url{https://www.omniglot.com/language/numbers/ughele.htm}
\end{itemize}
Frostad's 2012 grammar was cited by these sources, but no URL for it appeared in the results.
\section{Religious and Community Resources}
Joshua Project reports that Christianity is the main religion and that Scripture is available in Ughele. This was not verified: \url{https://joshuaproject.net/people_groups/15689/BP}
\section{Digital and Audio Resources}
\begin{itemize}
\item Verbix page on Ughele (content not checked): \url{https://docs.verbix.com/Languages/Ughele}
\end{itemize}
No audio was found.
\section{Curriculum Development Context}
Ughele is spoken by the whole community, but the only documentation found is a single grammar. No school programme or orthography standard was found. A curriculum would start from the Frostad grammar and from work with community members.
\end{document}